% lecture10.tex - 特征向量与特征值
% 第 10 章: 特征向量与特征值

\chapter{特征向量与特征值}
\label{ch:eigen}

\begin{abstract}
特征向量 (eigenvector) 是线性变换只拉伸而不改变其方向的特殊方向。这一概念是线性代数中最重要的思想之一, 在物理学、工程学和数据科学中有着广泛的应用。
\end{abstract}

% ============================================
\section{动机: 特殊的向量}
% ============================================

\begin{intuition}
当你施加一个线性变换时, 大多数向量都会偏离原来的方向。但有些特殊的向量只被拉伸或压缩, 仍然\textbf{保持在同一条直线上}。

这些特殊的向量称为\textbf{特征向量 (eigenvector)}。
\end{intuition}

\begin{center}
\begin{tikzpicture}[scale=1.2]
    % 变换前
    \begin{scope}[shift={(0,0)}]
        \draw[grid] (-2,-2) grid (2,2);
        \draw[axis] (-2,0) -- (2,0);
        \draw[axis] (0,-2) -- (0,2);
        \draw[vector, blue!70, very thick] (0,0) -- (1,0.5) node[right] {$\vv_1$};
        \draw[vector, red!70, very thick] (0,0) -- (0.5,1) node[above] {$\vv_2$};
        \draw[vector, green!60!black, very thick] (0,0) -- (-1,1) node[above left] {特征方向!};
        \node at (0,-2.5) {变换前};
    \end{scope}

    \draw[-{Stealth}, very thick] (2.5,0) -- (4,0) node[midway, above] {$T$};

    % 变换后
    \begin{scope}[shift={(6.5,0)}]
        \draw[blue!20] (-2,-1) -- (0,2) -- (4,2) -- (2,-1) -- cycle;
        \draw[axis] (-2,0) -- (4,0);
        \draw[axis] (0,-2) -- (0,2);
        \draw[vector, blue!70, very thick] (0,0) -- (1.5,1) node[right] {$T(\vv_1)$};
        \draw[vector, red!70, very thick] (0,0) -- (1.5,1.5) node[above] {$T(\vv_2)$};
        \draw[vector, green!60!black, very thick] (0,0) -- (-2,2) node[above left] {$2 \times$ 特征方向!};
        \node at (1,-2.5) {变换后};
    \end{scope}
\end{tikzpicture}
\end{center}

% ============================================
\section{正式定义}
% ============================================

\begin{definition}[特征向量与特征值]
非零向量 $\vv$ 是矩阵 $\mA$ 的\vocab{特征向量 (eigenvector)}, 如果:
\[
\mA\vv = \lambda\vv
\]
对某个标量 $\lambda$ 成立。该标量称为对应的\vocab{特征值 (eigenvalue)}。
\end{definition}

\begin{intuition}
$\mA\vv = \lambda\vv$ 的意思是:
\begin{center}
``变换 $\mA$ 作用在向量 $\vv$ 上, 就相当于一次标量乘法''
\end{center}
向量 $\vv$ 保持在同一条直线上, 只是被 $\lambda$ 缩放了。
\end{intuition}

\begin{example}
对于矩阵 $\mA = \twomat{3}{1}{0}{2}$:

向量 $\vv = \twovec{1}{0}$ 是一个特征向量:
\[
\mA\vv = \twomat{3}{1}{0}{2}\twovec{1}{0} = \twovec{3}{0} = 3\twovec{1}{0} = 3\vv
\]
对应的特征值为: $\lambda = 3$
\end{example}

% ============================================
\section{求特征值}
% ============================================

\begin{theorem}[特征方程]
$\lambda$ 是 $\mA$ 的特征值当且仅当:
\[
\det(\mA - \lambda\mI) = 0
\]
\end{theorem}

\begin{proof}
$\mA\vv = \lambda\vv$ 可写为:
\[
\mA\vv - \lambda\vv = \vzero \implies (\mA - \lambda\mI)\vv = \vzero
\]
该方程有非零解当且仅当 $\mA - \lambda\mI$ 不可逆, 即 $\det(\mA - \lambda\mI) = 0$。
\end{proof}

\begin{definition}[特征多项式]
$\det(\mA - \lambda\mI)$ 是关于 $\lambda$ 的多项式, 称为\vocab{特征多项式 (characteristic polynomial)}。它的根就是特征值。
\end{definition}

\begin{example}
对于 $\mA = \twomat{3}{1}{0}{2}$:
\[
\det(\mA - \lambda\mI) = \det\twomat{3-\lambda}{1}{0}{2-\lambda} = (3-\lambda)(2-\lambda) - 0 = 0
\]
根为: $\lambda_1 = 3$, $\lambda_2 = 2$
\end{example}

% ============================================
\section{求特征向量}
% ============================================

\begin{theorem}
对于每个特征值 $\lambda$, 对应的特征向量可通过求解齐次方程组得到:
\[
(\mA - \lambda\mI)\vv = \vzero
\]
\end{theorem}

\begin{example}
继续上一例, 取 $\lambda = 2$:
\[
(\mA - 2\mI)\vv = \twomat{1}{1}{0}{0}\twovec{v_1}{v_2} = \vzero
\]
方程为: $v_1 + v_2 = 0$, 于是 $v_1 = -v_2$

特征向量为: $\vv = t\twovec{-1}{1}$, 其中 $t \neq 0$ 为任意数
\end{example}

% ============================================
\section{特征空间}
% ============================================

\begin{definition}[特征空间]
对应于特征值 $\lambda$ 的\vocab{特征空间 (eigenspace)}:
\[
E_\lambda = \Null(\mA - \lambda\mI) = \{\vv \mid \mA\vv = \lambda\vv\}
\]
\end{definition}

\begin{intuition}
特征空间包含变换 $\mA$ 仅以倍数 $\lambda$ 缩放的所有向量。该空间总是一个向量\vocab{子空间 (subspace)}。
\end{intuition}

% ============================================
\section{几何解释}
% ============================================

\begin{practical}
\textbf{三维旋转}

对于 $\Rthree$ 中的旋转, 特征值 $\lambda = 1$ 对应的特征向量就是\textbf{旋转轴}! 该向量保持不动。

用轴和角描述旋转, 比用 $3 \times 3$ 矩阵简洁得多。
\end{practical}

\begin{example}
错切矩阵: $\mA = \twomat{1}{1}{0}{1}$

特征值: $\lambda = 1$ (重根)

特征向量: $\twovec{1}{0}$ (只有一个特征方向)

解释: 错切保持 $x$ 轴不动。
\end{example}

% ============================================
\section{特殊情形}
% ============================================

\begin{theorem}[二维旋转]
旋转矩阵 $\mR_\theta = \twomat{\cos\theta}{-\sin\theta}{\sin\theta}{\cos\theta}$, 当 $\theta \neq 0, \pi$ 时:

特征值为: $\lambda = \cos\theta \pm i\sin\theta$ (复数!)

没有实向量保持不动——所有向量都旋转了。
\end{theorem}

\begin{warning}
即使对实矩阵, 特征值也可能是\textbf{复数}! 旋转中就会出现这种情况。
\end{warning}

% ============================================
\section{应用}
% ============================================

\begin{practical}
\textbf{谷歌 PageRank}

把网页看作向量, 把链接看作矩阵。主特征向量 (对应最大特征值的特征向量) 反映了各网页的相对重要性。
\end{practical}

\begin{practical}
\textbf{主成分分析 (PCA)}

在机器学习中, 协方差矩阵的特征向量给出了数据变化的主要方向。
\end{practical}

\begin{practical}
\textbf{量子力学}

算子的特征值 = 可能的测量结果

特征向量 = 系统的稳定状态
\end{practical}

% ============================================
\section{练习}
% ============================================

\begin{exercise}
求下列矩阵的特征值与特征向量:
\[
\mA = \twomat{4}{1}{2}{3}
\]
\end{exercise}

\begin{exercise}
证明对角矩阵 (diagonal matrix) 的特征值就是它的对角元。
\end{exercise}

\begin{exercise}
求旋转 $90°$ 的旋转矩阵的特征值。
\end{exercise}

\begin{exercise}
若 $\lambda$ 是 $\mA$ 的特征值, 证明 $\lambda^2$ 是 $\mA^2$ 的特征值。
\end{exercise}

\begin{exercise}[挑战]
证明矩阵的迹 (对角元之和) 等于特征值之和。
\end{exercise}

\begin{problem}
种群矩阵:
\[
\mL = \twomat{0}{4}{0.5}{0}
\]
求主特征值并解释其含义。
\end{problem}
